% =======================================
% ПРЕАМБУЛА ДЛЯ КОНСПЕКТОВ ПО ИНФОРМАТИКЕ
% =======================================

% НАЧАЛО: Языки и шрифты
\usepackage{cmap}            % Поиск и копирование русского текста из PDF
\usepackage[utf8]{inputenc}  % Кодировка исходного текста
\usepackage[T2A]{fontenc}    % Поддержка кириллических шрифтов
\usepackage[russian]{babel}  % Правила русского языка и переносы строк

% ГЕОМЕТРИЯ: Размеры полей страницы
\usepackage[
    left=2cm,
    right=2cm,
    top=2cm,
    bottom=2cm
]{geometry}

% ГРАФИКА: Иллюстрации и рисунки
\usepackage{graphicx}        % Вставка готовых картинок (PNG, JPG, PDF)
\usepackage{tikz}            % Рисование графиков и схем кодом

% ТИПОГРАФИКА: Оформление текста
\usepackage{indentfirst}     % Принудительный отступ (красная строка) для всех абзацев
\usepackage{microtype}       % Выравнивание краев текста без разрывов слов

% НУМЕРАЦИЯ: Формат разделов
\renewcommand{\thesection}{§\,\arabic{section}} % Параграфы в виде: § 1

% ЛИСТИНГ: Форматирование исходного кода программ
\usepackage{verbatim}
\usepackage{multicol}

% КОНЕЦ: Интерактивность и ссылки
\usepackage[
    colorlinks=true,   % Цветной текст для ссылок вместо цветных рамок
    linkcolor=black,   % Черный цвет ссылок в оглавлении
    urlcolor=blue,     % Синий цвет для интернет-ссылок
    citecolor=black    % Черный цвет для библиографии
]{hyperref}
